\newpage
\pdfbookmark[1]{\infoname}{info} % Lisame infolehe PDF-faili järjehoidjatesse.

% Eesti keeles info.
\begin{info}
\begin{abstract}
Aines ``Automaadid, keeled ja translaatorid'' on \textsc{CMa} virtuaalmasina mudelit kasutatud juba mitu aastat kompilaatorite õpetamiseks, sest selle semantika sarnaneb paljudes aspektides reaalse täitmiskeskkonna tööpõhimõtetega. Bakalaureusetöö eesmärk oli laiendada \textsc{CMa} (\emph{C-Machine}) mudeli implementatsiooni, lisades olemasolevale koodibaasile funktsioonikutsete toe ja tagades, et keskkond jääb kogu käitamise vältel korrektseks. Rakendusliku uurimuse käigus analüüsiti olemasolevat implementatsiooni Javas, kavandati ja implementeeriti funktsioonikutsete mehhanism ning verifitseeriti keskkonna korrektsus, sealhulgas pesastatud ja rekursiivsete väljakutsete korral. Töö tulemusena lisati käsustikku 13 uut käsku, kirjutati ligikaudu 1251 rida uut koodi ning korrektsus verifitseeriti 46 JUnit testiga, mis katavad nii üksikute käskude kui ka pikemate programmide toimimise.
\end{abstract}

\keywords{abstraktne masin, virtuaalmasin, \textsc{CMa}, kompilaator, funktsioonikutse, Java}

\cercs{P175 Informaatika, süsteemiteooria}
\end{info}


% Inglise keeles info.
\begin{otherInfo}{english}{Extending the \textsc{CMa} Virtual Machine in Java to Support Function Calls}
\begin{abstract}
In the course ``Automata, languages and compilers'', the \textsc{CMa} virtual machine model has been used for several years in compiler education because its semantics resemble real execution environments in several important aspects. The aim of this bachelor's thesis was to extend the implementation of the \textsc{CMa} (\emph{C-Machine}) model by adding support for function calls while ensuring that the environment remained correct throughout execution. The applied research included an analysis of the existing implementation in Java, the design and implementation of a function call mechanism, and the verification of environment correctness, including nested and recursive calls. The implementation extended the instruction set with 13 new instructions, contributed approximately 1251 lines of new code, and correctness was verified by a suite of 46 JUnit tests covering the functionality of both individual instructions and complete programs.
\end{abstract}

\keywords{Abstract machine, virtual machine, \textsc{CMa}, compiler, function call, Java}

\cercs{P175 Informatics, systems theory}
\end{otherInfo}
